\section{Stage VI: The Outer Loop and a Register-Bound Kernel}
\label{sec:stage6}

\subsection{Invariant-Stride Induction Variables}
\label{sec:ivs}
With the inner loops nearly identical to C++, the difference moved to the
\emph{row-group} loop. GCC advances four row pointers with four \code{addq}s
per group; OxCaml rebuilt them with an \code{imulq} and eight \code{leaq}s.
The row base \code{o * cols} steps by \code{4 * cols}, an invariant but not
constant stride, which strength reduction declines:

\begin{lstlisting}[language=CmmX,caption={The row bases in Cmm (tagged arithmetic).},label=lst:rowbases,numbers=none]
b0 = (+ (* (+ o -1) (>>s cols 1)) 1)    ; a product with a runtime stride
b1 = (+ (+ b0 cols) -1)
b2 = (+ (+ b0 (<< cols 1)) -2)
b3 = (+ (+ b0 (* cols 3)) -3)
\end{lstlisting}

The invariant-stride pass generalises the affine model to
\code{pointer + (stride * counter + offset)}, with stride and offset
\emph{expressions} over variables bound outside the loop. The stride is
computed once and each back edge becomes one pointer addition; since the pass
carries the four interior \emph{pointers}, the \code{imulq}/\code{leaq} chain
becomes dead. Instruction selection now allows an \code{Addr} in a
loop-carried register under the pass's condition; the \code{bad GC root}
check still applies. The preamble fell from 25 to 13 instructions, the shape
GCC emits. Carrying the indices
instead (\code{-no-cmm-carried-pointers}) keeps nine values live instead of
four and spills in the hot loop: on the compilers built without polls,
carrying pointers is @P{ox54_ivs_noptr,ox54_ivs}@ faster, and the pass is
worth @P{ox54_ip,ox54_ivs}@ over interior pointers alone.

\subsection{Loop-Invariant Code Motion: A Designed Null Result}
The inner loop's trip count, \code{(cols - 2) >>s 1}, is recomputed for every
row group although it is invariant; flambda2 has no loop-invariant code
motion~\cite{aho2006dragon}, so the LICM pass hoists pure integer
\code{let}-bindings out of GC-free loops. Hoisting unconditionally made the
allocator spill a value inside the inner loop (visible in the assembly; the
measured $-1.1\%$, in a separate 15-round run, is within the noise): a hoist
is free only if its operands die, and \code{cols} does not. LICM therefore
uses strength reduction's register-neutrality condition and, on microgpt,
declines as designed: \code{pr_c_ivs} and \code{pr_counted} compile to the
same code, and the measured difference, @P{pr_c_ivs,pr_counted}@, shows the
run's resolution. Whether it hoists depends on the loop, not on this
benchmark.

\paragraph{Observation: spills track the result.} Every change that made the
allocator spill inside the inner loop lost (carried indices, unconditional
hoisting); every change that removed a live value or a spill won. Not every
spill is costly: an eight-row kernel (Section~\ref{sec:simd-why}) spills four
values per iteration yet runs as fast as the four-row kernel, because its L1
reloads, spread over eight multiply--adds, leave slightly fewer instructions
per multiply--add (about 3.1 against 3.25). A spill hurts when it adds work to
every multiply--add.

\subsection{The Ladder}
Table~\ref{tab:passes} gives each pass's contribution on the pull-request
compiler configured like the release. Counted loops (Section~\ref{sec:counted})
come first, so that polls sit at loop exits and every pass can act; source and
output are identical in every row.

\begin{table}[h]
  \centering
  \small
  \begin{tabularx}{\textwidth}{@{}lXrr@{}}
    \toprule
    step & what it does & paired gain (rounds) & gap \\
    \midrule
    passes off (\code{pr_off}) & polls on every back edge & & @G{pr_off}@ \\
    + counted loops & polls move to loop exits; nothing else changes & @P{pr_off,pr_c_off}@ & @G{pr_c_off}@ \\
    + strength reduction & carries \code{base + i} as a loop parameter & @P{pr_c_off,pr_c_sr}@ & @G{pr_c_sr}@ \\
    + interior pointers & one hoisted \code{Addr} per row, one shared index & @P{pr_c_sr,pr_c_ip}@ & @G{pr_c_ip}@ \\
    + invariant strides & carries the row pointers across the row loop & @P{pr_c_ip,pr_c_ivs}@ & @G{pr_c_ivs}@ \\
    + LICM (\code{pr_counted}) & declines here, by design (identical code) & @P{pr_c_ivs,pr_counted}@ & @G{pr_counted}@ \\
    \midrule
    all, from passes off & & \textbf{@D{pr_off,pr_counted}@} (@W{pr_off,pr_counted}@) & \\
    \bottomrule
  \end{tabularx}
  \caption{The pass ladder (pull-request compiler, unchanged source, run W);
  gap to scalar C++ in the last column.}
  \label{tab:passes}
\end{table}

Counted loops alone only remove polls; strength reduction and interior
pointers carry the gain. Invariant strides add nothing measurable here: with
counted loops, interior pointers already remove this kernel's row-group
arithmetic. With polls kept on every back edge, the four passes give only
@P{pr_off,pr_all}@ (Section~\ref{sec:reality}).

\subsection{Where Is the Rest of the Gap?}
\label{sec:residual}
With invariant strides the regions dominated by \code{linear} run at
$1.15$--$1.25\times$ of scalar C++ and attention, at $1.42\times$, is the
worst; with counted loops the picture is the same ($1.17$--$1.27\times$,
attention $1.49\times$; Figure~\ref{fig:regions}). In isolated calls
(Table~\ref{tab:shapes}, left), $16\times64$, with four times as many columns
per row, has a larger gap than $64\times16$ with equal work, so part of the
gap lives in the inner loop body.

\begin{table}[h]
  \centering
  \small
  \begin{tabular}{@{}lrrr@{\hspace{3em}}lr@{}}
    \toprule
    shape (rows $\times$ cols) & C++ ns & OxCaml ns & ratio & hand-written loop & ticks/iter \\
    \midrule
    $64\times16$ (fc1)      & 228.5 & 240.9 & $1.05\times$ & C++ shape & 2.875 \\
    $16\times16$ (q, k, v, o) & 57.4 & 62.0 & $1.08\times$ & OxCaml shape & 3.305 \\
    $16\times64$ (fc2)      & 209.5 & 250.0 & $1.19\times$ & \quad rotated, one branch & 3.301 \\
    $27\times16$ (lm\_head) & 89.0 & 111.6 & $1.25\times$ & \quad loads not folded & 3.301 \\
    $28\times16$ (no remainder) & 88.4 & 105.0 & $1.19\times$ & \quad scale 1 instead of 8 & 3.309 \\
    \bottomrule
  \end{tabular}
  \caption{Separate micro-benchmarks. Left: isolated \texttt{linear} calls,
  scalar C++ against the invariant-strides compiler. Right: hand-written 4-MAC
  loops, in time-stamp-counter ticks per iteration (3.0\,GHz counter, core at
  about 4.4\,GHz: one tick is about 1.46 core cycles).}
  \label{tab:shapes}
\end{table}

\paragraph{Loop structure.} Hand-written assembly versions of both loop shapes,
varied one difference at a time (Table~\ref{tab:shapes}, right), showed that
rotating the loop, un-folding the loads and changing the addressing scale are
each worth nothing. Only the encoding (two-operand SSE in C++, three-operand
VEX in OxCaml) moved the result, and not cleanly: run alone, the SSE loop
slowed from 2.875 to 3.675 ticks, an SSE/AVX transition effect. A cleaner
test: if SSE favoured C++, building C++ with VEX (\code{-mavx}) should slow it
towards OxCaml. In run W that build is @P{cpp_novec,cpp_novec_avx}@ against
the SSE build, so the SSE baseline is the harder one to match, not a
flattering one.

\paragraph{A model of the core.} Why is a loop with fewer instructions than
C++'s slower? A loop can be limited by instruction issue, by its execution
ports, or by chains of dependent instructions; only the first depends directly
on instruction count. \code{llvm-mca} with its Alder Lake description (the closest
available model of this machine's Meteor Lake performance cores) predicts which
limit applies
(Table~\ref{tab:mca}; a floating-point addition takes 2 cycles and a
multiplication 3). It puts the patched 13-instruction loop level with C++'s
15: both push the same eight floating-point operations through the same two
ports, at least three cycles per iteration, so, in the model, removing integer
instructions cannot make it faster than C++'s. The 5.2.0+ox loop's 23
instructions exceed what the core can issue in three cycles, which is why
removing instructions paid off until the loop reached C++'s size.

\begin{table}[h]
  \centering
  \small
  \begin{tabular}{@{}lrrrl@{}}
    \toprule
    inner loop & instructions & predicted & measured & predicted bound \\
    \midrule
    C++ scalar (\code{cpp_novec}) & 15 & 3.02 & 3.22 & floating-point ports \\
    OxCaml, counted loops (\code{pr_counted}) & 13 & 3.02 & 3.66 & floating-point ports \\
    OxCaml, SSE2 kernel, no polls & 16 & 3.02 & 3.46 & shuffle port \\
    OxCaml 5.2.0+ox release & 23 & 4.02 & 5.84 & instruction issue \\
    \bottomrule
  \end{tabular}
  \caption{Core cycles per 4-MAC iteration: predicted by \code{llvm-mca}, and
  measured with each loop copied at its original alignment into a timing
  harness (a separate measurement, 256 columns).}
  \label{tab:mca}
\end{table}

Measured alone, the patched loop is about 14\% slower than C++'s, where the
model predicts no difference; the cause is not identified here. This bounds
the rest of the gap: if the inner loop were the whole story, the program would
be at most about 14\% slower than C++; it is @G{pr_counted}@ behind, so
by this estimate roughly half of the remaining gap lies outside the inner
loop, in the
row-group setup (once per four rows of only 16 columns) and in the other
regions, attention above all.

\paragraph{Attention.} The attention loops call \code{linear} and
\code{softmax}, so the pointer passes do not apply. Swapping the value
gather's loop nest to do four MACs per iteration over contiguous memory took
OxCaml's attention block from 395 to 246\,ns in an isolated test, level with
unchanged C++ (248\,ns), which would gain too.
